% ============================================================
% Recreational Marijuana Legalization and Drug-Overdose
% Mortality: A Cross-Border Distance Design
% Formatted for ECONOMICS LETTERS (Elsevier, elsarticle class).
% REVISION (2026-07-25): full rewrite onto the stacked multi-
% period panel (every state's actual legalization date, 2020-
% 2025), replacing the earlier two-snapshot (Dec 2020/Dec 2023)
% design. Incorporates: a formal pre-trends check and
% mean-reversion correction, a medical-legalization placebo
% check, a persistence (post-switch dynamics) check, a county
% fixed-effects robustness check, and a corrected welfare
% translation using (life expectancy - average age at death)
% rather than full remaining actuarial life expectancy.
% NOTE: Elsevier submissions typically require Highlights as a
% SEPARATE uploaded file; included in-document here for a
% self-contained draft.
% Compile: pdfLaTeX -> BibTeX -> pdfLaTeX -> pdfLaTeX
% ============================================================
\documentclass[preprint,12pt,authoryear]{elsarticle}

\usepackage{amsmath,amssymb}
\usepackage{graphicx}
\usepackage{booktabs}
\usepackage{enumitem}
\usepackage[hidelinks]{hyperref}

\journal{Economics Letters}

\begin{document}

\begin{frontmatter}

\title{Recreational Marijuana Legalization and Drug-Overdose Mortality:\\
A Cross-Border Distance Design}

\author{Author name(s) withheld for review}

\begin{abstract}
\noindent We test whether proximity to a legal recreational-marijuana
market reduces drug-overdose mortality, using a stacked multi-period
panel of counties in 31 never-legalizing U.S. states and every state's
actual legalization date, 2020--2025, rather than a single
before/after snapshot. Switcher counties (distance to the nearest
legal-market county crossing below 50 miles) see overdose deaths fall
2.32 more per 100,000 than controls in the same transition
($p=0.004$); a continuous per-mile specification agrees in sign and is
itself significant ($p<0.001$). A formal pre-trends test shows the raw
switcher comparison is contaminated by mean reversion, which a
baseline-rate control corrects; the continuous specification, unlike
the switcher indicator, remains significant even under county fixed
effects, the most demanding test we run. A placebo check using
medical- rather than recreational-marijuana timing is null. The effect
does not compound in years after switching. Applying the national age
distribution of overdose deaths, the switcher effect implies roughly
755 averted deaths and 22,500 years of life saved -- illustrative, not
confirmed.
\end{abstract}

\begin{keyword}
recreational marijuana \sep opioid overdose mortality \sep cross-border spillover \sep staggered difference-in-differences \sep JEL: I18, I12, K42
\end{keyword}

\end{frontmatter}

\vspace{0.3cm}
\noindent\textbf{Highlights}
\begin{itemize}[leftmargin=1.4em]
\item Stacked multi-period cross-border distance design, using every state's actual legalization date, 2020--2025.
\item Switcher counties show 2.32 fewer deaths/100k ($p{=}0.004$); the continuous per-mile result agrees in sign ($p{<}0.001$).
\item A formal pre-trends test finds mean reversion in the raw switcher comparison; a baseline-rate control fixes it.
\item The continuous result, unlike the switcher result, survives county fixed effects -- the most demanding robustness test.
\item A medical-legalization placebo check is null; the effect does not compound in years after switching.
\item Implies $\sim$755 averted deaths, $\sim$22,500 life-years saved -- illustrative, not confirmed.
\end{itemize}

\vspace{0.4cm}

% ============================================================
% 1. INTRODUCTION
% ============================================================
\section{Introduction}

Does proximity to a legal recreational-marijuana market change
drug-overdose mortality for residents who never cross the border
themselves? The evidence on marijuana-opioid substitution is
unsettled: \citet{Bachhuber2014} found state medical-marijuana laws
associated with fewer opioid deaths, \citet{Shover2019} showed the
association reverses sign with more years of data, and
\citet{PowellPaculaJacobson2018} found any effect specific to
dispensary-based designs. \citet{MathurRuhm2022} conclude retail
dispensaries are, if anything, associated with \emph{higher} opioid
mortality. Most recently, \citet{Denkyirah2026} use a state-level
two-way-fixed-effects design for \emph{recreational} legalization and
find a decline of roughly 3.5 deaths per 100,000 -- a magnitude we
return to below.

Every one of these papers identifies off policy \emph{timing}. None
uses \emph{geography}: whether proximity to an external legal market
predicts a non-adopting county's own overdose trajectory, in the
spirit of \citet{HaoCowan2020}'s cross-border design for
marijuana-possession arrests. We apply that design to overdose
mortality using both a binary switcher indicator and a continuous,
per-mile distance measure, built as a stacked multi-period panel that
uses all 24 states' actual legalization dates rather than collapsing a
multi-year staggered rollout into two snapshots. An earlier,
two-snapshot version of this design found a null, sign-flipped
continuous estimate that we could not resolve with a bandwidth-based
fix; the multi-period panel resolves it directly. The multi-period
panel also makes possible three checks unavailable in the two-snapshot
version: a formal pre-trends test, a persistence (post-switch
dynamics) check, and a county fixed-effects robustness test -- all
reported below, together with a placebo-style check that swaps in
medical- for recreational-marijuana legalization timing.

% ============================================================
% 2. DATA AND EMPIRICAL STRATEGY
% ============================================================
\section{Data and empirical strategy}

We use 2,331 counties in the 31 U.S. states with no
recreational-marijuana retail market as of our latest
snapshot.\footnote{The companion project's original 32-state
never-legalizing panel included Maine; Maine's retail market opened
October 2020, before even our earliest snapshot, so its counties are
excluded here.} Each county's distance to the nearest legal-market
county is computed via Haversine distance between county centroids
(Census Gazetteer), with ``legal-market'' defined by first retail
sales date, for every December 2020--2025 -- the full window CDC
publishes county-level 12-month-ending overdose data for. Rather than
compare two fixed snapshots, we use all 24 states' actual
retail-opening dates and build a stacked panel of year-over-year
transitions. In transition $t-1\to t$, a county is a ``new switcher''
if its distance to the nearest legal-market county crossed from at
least 50 miles to below 50; a ``control'' if it stayed at least 50
miles away at both endpoints. Once a county switches it is dropped
from every later transition, avoiding the already-treated-unit-as-
control bias documented in the staggered difference-in-differences
literature \citep{GoodmanBacon2021}. Counties within 50 miles as of
the December-2020 baseline (116 of 2,331) are ``always-near'' and
dropped throughout.

The unit of observation is a county-transition pair $(i,t)$: county
$i$ contributes one row for each one-year transition
$t\in\{\text{2020--21},\ldots,\text{2024--25}\}$ it is ``at risk'' of
switching in -- every transition up to and including the one in which
it switches, if it ever does, after which it is dropped, as just
described. $\Delta Y_{it}$ is the change in county $i$'s overdose
death rate per 100,000 between the start and end of transition $t$.
$\text{Switcher}_{it}$ is a binary indicator equal to 1 if and only if
county $i$'s distance to the nearest legal-market county crossed from
at least 50 miles to below 50 miles \emph{during} transition $t$, so a
given county contributes at most one $\text{Switcher}_{it}=1$
observation across the whole panel; it equals 0 for every transition a
county has not yet switched in, including every transition for
counties that never switch. $\Delta\text{Dist}_{it}$ is the raw change
in miles to the nearest legal-market county over transition $t$,
defined for every county-transition row regardless of switcher status
(control counties can still move closer or farther without crossing
the threshold). $\delta_t$ denotes transition (period) fixed effects,
common to all counties observed in a given transition. Standard errors
are clustered by state, the level at which most of the underlying
legal-market-timing variation occurs. We pool all (county, transition)
observations into one stacked regression
\citep{CengizDubeLindnerZipperer2019}:
\begin{equation}
  \Delta Y_{it} = \alpha + \gamma \cdot \text{Switcher}_{it} + \delta_t + \varepsilon_{it}, \qquad
  \Delta Y_{it} = \alpha + \beta \cdot \Delta\text{Dist}_{it} + \delta_t + \varepsilon_{it}.
  \label{eq:main}
\end{equation}
Equation~\ref{eq:main} imposes one switcher effect $\gamma$ and one
distance effect $\beta$ common to all five transitions. We also
estimate a more general version that lets both vary freely by period,
via a full interaction with the transition dummies:
\begin{equation}
  \Delta Y_{it} = \alpha + \sum_t \gamma_t \cdot (\text{Switcher}_{it} \times \mathbb{1}[t])
                        + \sum_t \beta_t \cdot (\Delta\text{Dist}_{it} \times \mathbb{1}[t])
                        + \delta_t + \varepsilon_{it},
  \label{eq:period}
\end{equation}
which nests Equation~\ref{eq:main} as the special case
$\gamma_t=\gamma$ and $\beta_t=\beta$ for all $t$;
Section~\ref{sec:periodeffects} reports the period-specific
coefficients and tests whether pooling is appropriate.

As of December 2020, ten states had an open retail market; by December
2025, 23 did (Figure~\ref{fig:rollout}). Virginia has never opened a
general retail market and so never enters the legal-market set.
County-level overdose mortality comes from CDC's Vital Statistics
Rapid Release (VSRR) provisional county data (data.cdc.gov); CDC
WONDER is not reachable from our environment. NCHS confidentiality
rules suppress county-transition cells with 1--9 deaths at either
endpoint, leaving 4,259 of 10,313 potential observations (177 switch
events).

As a placebo-style check, we rebuild the identical stacked design
using each state's medical-marijuana access date (first regulated
medical dispensary sales, compiled for 39 states) in place of the
recreational-retail date, restricting the county universe to the 17 of
31 states not yet medically legal as of the December-2020 baseline.

\begin{figure}[htbp]
  \centering
  \includegraphics[width=\textwidth]{figures/fig1_rollout_maps_multiperiod.png}
  \caption{Legal-market rollout (by first retail sales date) and 50-mile treatment frontier, 2020--2025.}
  \label{fig:rollout}
\end{figure}

% ============================================================
% 3. RESULTS
% ============================================================
\section{Results}

\begin{table}[htbp]
\centering
\caption{Effect on drug-overdose mortality, preferred specification}
\label{tab:main}
\begin{tabular}{@{}lcc@{}}
\toprule
 & Switcher (binary) & Distance, per mile \\
\midrule
Overdose rate per 100k & $-2.320^{***}$ (0.800) & $+0.0122^{***}$ (0.0034) \\
 & $p=0.004$ & $p<0.001$ \\
\bottomrule
\end{tabular}
\vspace{2pt}
\begin{minipage}{0.85\textwidth}
\footnotesize\emph{Notes:} Outcome is the transition-level change in
overdose deaths per 100{,}000. Both columns include transition fixed
effects and a control for each county's baseline (2020) overdose-rate
level (Section~\ref{sec:pretrends}). SE clustered by state. $N=3{,}794$.
The distance coefficient is signed as the effect of an
\emph{increase} in distance, so it agrees in direction with the
(negative) switcher coefficient. $^{***}p<0.01$.
\end{minipage}
\end{table}

\begin{figure}[htbp]
  \centering
  \includegraphics[width=0.85\textwidth]{figures/fig2_overdose_coefplot_multiperiod.png}
  \caption{Effect on overdose mortality, recreational (primary) vs. medical (placebo) design (95\% CI).}
  \label{fig:coef}
\end{figure}

Switcher counties see overdose deaths fall 2.32 more per 100,000 than
controls in the same transition ($p=0.004$; Table~\ref{tab:main}) --
somewhat below \citet{Denkyirah2026}'s state-level estimate of roughly
3.5, though the same direction and order of magnitude. The continuous
per-mile specification agrees in direction and is itself significant
($p<0.001$).

\subsection{Pre-trends check and robustness to fixed effects}
\label{sec:pretrends}

Because eventual-switcher counties contribute observations to every
transition before they switch, we can test whether their overdose
trend was already diverging from never-switching counties
\emph{before} switching. The raw version of this test fails: leads on
event-time indicators (one and two-plus transitions before switching)
are large and jointly significant ($p<0.001$), comparable in magnitude
to the switch-transition effect itself, even though transition fixed
effects are already present. We traced this to mean reversion --
counties with unusually high 2020 overdose rates mechanically decline
more -- and adding each county's baseline (2020) rate as a control
resolves it (joint F-test on both leads, $p=0.318$; Figure~\ref{fig:pretrends}).
We also checked Census-division-by-transition fixed effects (a
geographic-clustering explanation); this also passes the joint test on
this panel but leaves the switch coefficient less precisely estimated
and proved sensitive to exact sample construction across our
diagnostic checks, so we prefer the baseline-rate control
(Table~\ref{tab:main}).

\begin{figure}[htbp]
  \centering
  \includegraphics[width=0.85\textwidth]{figures/fig5_eventstudy_pretrends.png}
  \caption{Pre-trends check: raw vs. mean-reversion-adjusted event-study coefficients.}
  \label{fig:pretrends}
\end{figure}

\begin{table}[htbp]
\centering
\caption{Robustness of main coefficients to fixed-effects controls}
\label{tab:fe}
\begin{tabular}{@{}lcccc@{}}
\toprule
 & \multicolumn{2}{c}{Switcher (binary)} & \multicolumn{2}{c}{Distance (per mile)} \\
Specification & Coef. & $p$ & Coef. & $p$ \\
\midrule
No controls & $-3.151$ & 0.005 & $+0.0128$ & $<0.001$ \\
+ Baseline overdose-rate level (preferred) & $-2.320^{***}$ & 0.004 & $+0.0122$ & $<0.001$ \\
+ State fixed effects & $-1.411$ & 0.144 & $+0.0129$ & $<0.001$ \\
+ County fixed effects & $-0.999$ & 0.532 & $+0.0122^{**}$ & 0.023 \\
\bottomrule
\end{tabular}
\vspace{2pt}
\begin{minipage}{0.92\textwidth}
\footnotesize\emph{Notes:} Each row is a separate regression on the
same $N=3{,}794$ sample; SE clustered by state. County fixed effects
use only within-county variation across a given county's (at most
five) transition-rows, the most demanding test we run -- 225 of 1,249
counties in this sample have only one transition-row and so contribute
no within-county variation. $^{**}p<0.05$, $^{***}p<0.01$.
\end{minipage}
\end{table}

Table~\ref{tab:fe} shows the two specifications respond very
differently to fixed-effects controls. The continuous per-mile
coefficient is essentially unmoved by state \emph{or} county fixed
effects and remains significant throughout. The binary switcher
coefficient survives the baseline-rate control but loses significance
under state fixed effects and further under county fixed effects,
indicating a meaningful share of its identifying variation comes from
\emph{which} counties or states happen to border a soon-to-open
market, not from finer within-unit proximity differences. We therefore
treat the continuous per-mile result as this paper's primary
quantitative claim, with the switcher result as corroborating but
comparatively more fragile.

\subsection{Period-specific effects}
\label{sec:periodeffects}

Equation~\ref{eq:main} assumes the switcher and distance effects are
constant across all five transitions. Estimating Equation~\ref{eq:period}
instead -- letting both vary freely by period -- and restricting to
the three transitions with enough switch events to identify a
period-specific coefficient (2020--2021 and 2024--2025 have zero
switch events between them after the baseline-rate restriction, so
contribute to the period fixed effects but not to their own
switcher/distance coefficient), gives Table~\ref{tab:period}.

\begin{table}[htbp]
\centering
\caption{Period-specific switcher and distance coefficients}
\label{tab:period}
\begin{tabular}{@{}lcccc@{}}
\toprule
 & \multicolumn{2}{c}{Switcher (binary)} & \multicolumn{2}{c}{Distance (per mile)} \\
Transition & Coef. & $p$ & Coef. & $p$ \\
\midrule
2021--2022 & $-0.648$ & 0.547 & $+0.0014$ & 0.717 \\
2022--2023 & $-1.718$ & 0.449 & $-0.0027$ & 0.802 \\
2023--2024 & $+8.280^{***}$ & 0.008 & $+0.1132^{***}$ & $<0.001$ \\
\bottomrule
\end{tabular}
\vspace{2pt}
\begin{minipage}{0.85\textwidth}
\footnotesize\emph{Notes:} One joint regression, all period
interactions plus period fixed effects and the baseline-rate control;
$N=3{,}794$; SE clustered by state. An F-test rejects that the three
switcher coefficients are equal ($p=0.002$) and that the three
distance coefficients are equal ($p=0.003$): the pooled
Table~\ref{tab:main} coefficient is not simply the same effect
repeated five times. $^{***}p<0.01$.
\end{minipage}
\end{table}

The 2023--2024 transition -- which contributes 75 of the panel's 177
switch events, concentrated in Ohio-border states (Ohio, West
Virginia, Pennsylvania, Kentucky, Indiana) -- has a positive,
significant coefficient on both specifications, the opposite sign from
the pooled result. We investigated this rather than reporting it at
face value, and three things resolve it. First, in raw, unadjusted
terms switcher counties in this transition declined \emph{more} than
controls (mean $\Delta$ overdose rate $-14.5$ vs.\ $-8.5$); the
positive sign appears only after the baseline-rate control, because
switchers in this specific transition have roughly double the 2020
baseline rate of controls (53.4 vs.\ 26.1), a far more severe imbalance
than in other transitions. Second, this period's control pool is
geographically mismatched to its switcher pool -- comparing
Appalachian counties against a national control pool concentrated in
the South Atlantic and West South Central divisions -- and adding
Census-division fixed effects shrinks the switcher coefficient to
$+2.79$ ($p=0.253$) and the distance coefficient to $+0.015$
($p=0.335$), both no longer significant. Third, a handful of
extreme-outlier small counties (implausibly large overdose-rate swings
driven by small population denominators) have outsized leverage on
this single-period fit; excluding the top 5\% by Cook's distance
likewise removes significance ($+1.82$, $p=0.166$). Combining division
fixed effects with winsorizing the baseline-rate control at its 95th
percentile gives $+2.33$ ($p=0.304$) for the switcher coefficient and a
smaller, though still marginally significant, $+0.032$ ($p=0.036$) for
distance. We read this as within-period estimation noise from
isolating a geographically concentrated, high-baseline-rate cohort with
few observations, not as evidence against the pooled result: it draws
on switch events spread across many more states and periods than any
single transition, and Table~\ref{tab:fe}'s state and county
fixed-effects checks already show the continuous specification is not
being driven by any one state or region.

\subsection{Medical-legalization placebo check}

Re-running the identical stacked design with medical- in place of
recreational-marijuana timing, both the switcher ($b=+0.51$,
se$=1.04$, $p=0.62$) and continuous ($b=-0.015$, se$=0.008$, $p=0.08$)
estimates are null -- consistent with the recreational-retail effect
reflecting retail access specifically rather than the design
mechanically detecting any nearby marijuana-policy change
(Figure~\ref{fig:coef}, bottom two rows).

\subsection{Persistence check}

Does the switcher effect continue, flatten, or reverse in the years
after a county crosses the 50-mile threshold? We re-estimate
Equation~\ref{eq:main} with event-time indicators for the switch
transition and each subsequent transition (one, and two-or-more,
transitions after switching), keeping switched counties in the panel
afterward purely for this check (they are otherwise dropped, per
Section~2). With the baseline-rate control, the switch-transition
effect remains significant ($-2.42$, $p=0.001$), but neither
subsequent-transition coefficient is significant (one transition
after: $+1.22$, $p=0.171$; two-plus: $+1.14$, $p=0.157$;
Figure~\ref{fig:persistence}), and a joint test rejects that the lag
effects equal the switch effect ($p=0.006$). The one-year decline
associated with switching does not continue to compound in later
years; if anything, the point estimates hint at partial reversion,
though not significantly so.

\begin{figure}[htbp]
  \centering
  \includegraphics[width=0.85\textwidth]{figures/fig6_persistence.png}
  \caption{Persistence check: switch-transition effect vs. subsequent-transition effects.}
  \label{fig:persistence}
\end{figure}

\subsection{Welfare translation}

National age-specific overdose rates (CDC NCHS Data Query System,
2024) imply a population-weighted average age at overdose death of
46.1 years. We translate the switcher coefficient into years of life
saved using the gap between life expectancy and this average age at
death, rather than the full remaining actuarial life expectancy, since
a person who dies of an overdose at 46 does not have the population's
full remaining life expectancy at stake, only the gap between when
they died and how long they could plausibly have lived absent this
specific cause. Life expectancy is drawn from CDC's state
life-expectancy-at-birth estimates \citep{CDCStateLE2021},
population-weighted across the 177 switcher counties' states (75.9
years), for a gap of 29.7 years per averted death. Applying the
preferred switcher coefficient to the 32.5 million residents of the
177 non-suppressed switcher counties implies approximately 755 averted
deaths (95\% CI: 245 to 1{,}265) over 2020--2025, and roughly 22,500
years of life saved (95\% CI: 7,300 to 37,600;
Figure~\ref{fig:agedist}). This is illustrative of what the point
estimate would mean, not a confirmed welfare effect: it inherits every
caveat on the underlying coefficient, including its greater
sensitivity to fixed-effects controls (Section~\ref{sec:pretrends}).

\begin{figure}[htbp]
  \centering
  \includegraphics[width=\textwidth]{figures/figC_age_distribution.png}
  \caption{National age distribution of drug-overdose deaths, 2024.}
  \label{fig:agedist}
\end{figure}

% ============================================================
% 4. CONCLUSION
% ============================================================
\section{Conclusion}

Extending the cross-border distance design to a stacked multi-period
panel that uses every state's actual legalization date resolves a
fragility in an earlier two-snapshot version: the continuous per-mile
specification is now sign-consistent with, and more robust than, the
binary switcher result, surviving even county fixed effects. A formal
pre-trends test finds and corrects a mean-reversion problem in the raw
switcher comparison; a persistence check finds the switch-year effect
does not compound in subsequent years; and a medical-legalization
placebo check is null, as it should be if the effect is specific to
retail-market access. Main limitations: VSRR data are provisional and
small-cell suppression removes more than half the panel, concentrated
where identifying variation is richest; the 50-mile radius is a design
choice; the panel is bounded to 2020--2025; and the welfare translation
mixes a state-level life-expectancy figure with a national age-at-death
figure. We recommend this be treated as suggestive evidence of an
opioid-substitution channel specific to recreational retail-market
access, meriting confirmation with final, non-provisional NCHS
mortality data.

% ============================================================
% SUPPLEMENTARY MATERIAL
% ============================================================
\appendix
\section{Supplementary figures}
\label{sec:appendix}

\noindent\emph{Full robustness-check summary.} Table~\ref{tab:robustsummary}
consolidates every fixed-effects and sample-construction check reported
in this paper into one place. Panel A reproduces Table~\ref{tab:fe}
(the checks that bear directly on the pooled, headline specification
in Table~\ref{tab:main}). Panel B reports the additional checks used
only to diagnose and resolve the 2023--2024 period-specific sign
reversal discussed in Section~\ref{sec:periodeffects}; these are
estimated on that single transition in isolation ($n=691$) and were
never substituted into the pooled specification.

\begin{table}[htbp]
\centering
\caption{Full robustness-check summary}
\label{tab:robustsummary}
\begin{tabular}{@{}lcccc@{}}
\toprule
 & \multicolumn{2}{c}{Switcher (binary)} & \multicolumn{2}{c}{Distance (per mile)} \\
Specification & Coef. & $p$ & Coef. & $p$ \\
\midrule
\multicolumn{5}{@{}l}{\emph{Panel A: Pooled specification, all periods ($N=3{,}794$) -- bears on Table~\ref{tab:main}}} \\
No controls & $-3.151$ & 0.005 & $+0.0128^{***}$ & $<0.001$ \\
+ Baseline overdose-rate level (preferred) & $-2.320^{***}$ & 0.004 & $+0.0122^{***}$ & $<0.001$ \\
+ State fixed effects & $-1.411$ & 0.144 & $+0.0129^{***}$ & $<0.001$ \\
+ County fixed effects & $-0.999$ & 0.532 & $+0.0122^{**}$ & 0.023 \\
\addlinespace
\multicolumn{5}{@{}l}{\emph{Panel B: 2023--2024 transition only ($n=691$) -- diagnoses Section~\ref{sec:periodeffects}'s sign reversal}} \\
Baseline-rate control only & $+8.280^{***}$ & 0.008 & $+0.1132^{***}$ & $<0.001$ \\
+ Census-division fixed effects & $+2.786$ & 0.253 & $+0.0149$ & 0.335 \\
+ Winsorized baseline rate (95th pctile) & $+4.574^{**}$ & 0.032 & --- & --- \\
+ Division FE \& winsorized baseline & $+2.328$ & 0.304 & $+0.0320^{**}$ & 0.036 \\
Excluding top 5\% by Cook's distance & $+1.822$ & 0.166 & --- & --- \\
\bottomrule
\end{tabular}
\vspace{2pt}
\begin{minipage}{0.95\textwidth}
\footnotesize\emph{Notes:} SE clustered by state throughout. Panel A:
each row a separate regression on the identical $N=3{,}794$ pooled
sample. Panel B: each row a separate regression restricted to the
2023--2024 transition only; ``---'' denotes specifications where the
distance coefficient was not separately estimated. Panel B's raw,
unadjusted group means (no regression) show switcher counties declined
\emph{more} than controls in this transition ($-14.5$ vs.\ $-8.5$
deaths per 100,000); the positive sign in Panel B's first row appears
only once a linear baseline-rate control is applied to a period with
an unusually large baseline-rate imbalance between switchers and
controls (Section~\ref{sec:periodeffects}). $^{**}p<0.05$,
$^{***}p<0.01$.
\end{minipage}
\end{table}

\noindent\emph{Nonlinearity check.} Table~\ref{tab:main}'s continuous
specification imposes a constant marginal effect. Adding a quadratic
term in $\Delta\text{Dist}_{it}$, the quadratic term is itself
significant ($p=0.005$; F-test, $p=0.009$; joint F-test with a cubic
term, $p<0.001$) -- unlike an earlier two-snapshot version of this
design, where no nonlinearity was detectable.
Figure~\ref{fig:doseresponse} plots the dose-response relationship and
implied marginal effect.

\begin{figure}[htbp]
  \centering
  \includegraphics[width=\textwidth]{figures/figA_doseresponse_multiperiod.png}
  \caption{Dose-response (panel a) and implied marginal effect (panel b), pooled stacked-transition panel.}
  \label{fig:doseresponse}
\end{figure}

\noindent\emph{Diff-in-discontinuities check.} We also tested for a
discontinuous jump in overdose mortality exactly at the 50-mile
threshold, with the bandwidth on each side chosen by the CCT
MSE-optimal selector \citep{CalonicoCattaneoTitiunik2014} in each
year's cross-section separately, differencing the 2025 and 2020
estimates following \citet{GrembiNanniciniTroiano2016}. This check is
inconclusive: the only significant year is 2020 (the baseline year,
before almost any rollout had occurred, so not itself interpretable as
a treatment effect), and the discontinuity shrinks rather than
strengthens by 2025 (diff-in-disc $=-14.68$, se$=8.79$, $p=0.095$). We
read this as a reason for some caution about the near/far comparison
generally, not as confirming or overturning the main results.

\noindent\emph{State-level trends and legalization timing.}
Figure~\ref{fig:statetrends} plots monthly overdose mortality,
2015--2026, for all 24 states in the source timing table, with both
the policy-effective date and the first-retail-sales date marked.
Figure~\ref{fig:eventstudy} pools the 23 states with a retail date,
aligned by months since first retail sales; neither figure feeds into
the county-level regression design.

\begin{figure}[htbp]
  \centering
  \includegraphics[width=\textwidth]{figures/figD_state_overdose_trends_new.png}
  \caption{State overdose mortality trends, 2015--2026, with policy-effective and retail-sales dates.}
  \label{fig:statetrends}
\end{figure}

\begin{figure}[htbp]
  \centering
  \includegraphics[width=\textwidth]{figures/figE_eventstudy_overdose_new.png}
  \caption{Overdose mortality around the retail-legalization cutoff, pooled across 23 states.}
  \label{fig:eventstudy}
\end{figure}

% ============================================================
% REFERENCES
% ============================================================
\bibliographystyle{apalike}
\bibliography{references}

\end{document}
